%%%%%%%%%%%%%%%%%%%%%%%%%%%%%%%%%%%%%%%%%
% "ModernCV" CV and Cover Letter
% LaTeX Template
% Version 1.1 (9/12/12)
%
% This template has been downloaded from:
% http://www.LaTeXTemplates.com
%
% Original author:
% Xavier Danaux (xdanaux@gmail.com)
%
% License:
% CC BY-NC-SA 3.0 (http://creativecommons.org/licenses/by-nc-sa/3.0/)
%
% Important note:
% This template requires the moderncv.cls and .sty files to be in the same
% directory as this .tex file. These files provide the resume style and themes
% used for structuring the document.
% latexmk -pdv -pvc main.tex
%%%%%%%%%%%%%%%%%%%%%%%%%%%%%%%%%%%%%%%%%

%----------------------------------------------------------------------------------------
%	PACKAGES AND OTHER DOCUMENT CONFIGURATIONS
%----------------------------------------------------------------------------------------

\documentclass[11pt,a4paper,sans]{moderncv} % Font sizes: 10, 11, or 12; paper sizes: a4paper, letterpaper, a5paper, legalpaper, executivepaper or landscape; font families: sans or roman
\usepackage{standalone}
\moderncvstyle{classic} % CV theme - options include: 'casual' (default), 'classic', 'oldstyle' and 'banking'
\moderncvcolor{grey} % CV color - options include: 'blue' (default), 'orange', 'green', 'red', 'purple', 'grey' and 'black'


\usepackage[scale=0.85]{geometry} % Reduce document margins
%\setlength{\hintscolumnwidth}{3cm} % Uncomment to change the width of the dates column
%\setlength{\makecvtitlenamewidth}{10cm} % For the 'classic' style, uncomment to adjust the width of the space allocated to your name

%\usepackage[utf8]{inputenc}

%\usepackage{booktabs}
\usepackage{fontawesome}
\usepackage{tikz}%
\usepackage{marvosym} % For cool symbols.
\usepackage{lastpage}
\usepackage{verbatim} % for comment block
\rfoot{\addressfont\itshape\textcolor{gray}{Page \thepage\ of \pageref{LastPage}}}


\tikzset{curve in scope/.style={
		every path/.style={
			draw=white,
			double distance=1pt,
			line width=2pt,
			double=black,
		},
		every node/.style={
			color=blue
		}
	}
}

%----------------------------------------------------------------------------------------
%	NAME AND CONTACT INFORMATION SECTION
%----------------------------------------------------------------------------------------

\firstname{Kiran} % Your first name
\familyname{Makhijani} % Your last name

% All information in this block is optional, comment out any lines you don't need
%\title{Curriculum Vitae}
%\address{Bay Area CA 95030}
\mobile{(+1) 408-410-2484}


%\fax{(000) 111 1113}

%\social{github}{stefano-bragaglia}
\email{kiran@mirani.net}


%\homepage{www.github.io/}kiranmak}

% social link \faGithub, \faSkype, \faLinkedin,\faStackExchange, and \faStackOverflow
\extrainfo{
    \faGithub\href{https://github.com/kiranmak}{ Github} \quad
    \faLinkedin\href{https://www.linkedin.com/kiranmak/}{ Linkedin} \quad
    %\faSkype\href{https://skype.com/abc}{Skype}
    }



%\social[linkedin][www.linkedin.com]{name}
% The first argument is %the url for the clickable link, the second argument is the url displayed in the %template - this allows special characters to be displayed such as the tilde in this %example

%\photo[40pt][0.2pt]{picture} % The first bracket is the picture height, the second is %the thickness of the frame around the picture (0pt for no frame)
%\quote{Not Attention, Patience is all we need.}

%----------------------------------------------------------------------------------------

\newcommand{\cvdoublecolumn}[2]{%
  \cvitem[.75em]{}{%
    \begin{minipage}[t]{\listdoubleitemcolumnwidth}#1\end{minipage}%
    \hfill%
    \begin{minipage}[t]{\listdoubleitemcolumnwidth}#2\end{minipage}%
    }%
}

\usepackage{multibbl}
\newcommand\Colorhref[3][orange]{\href{#2}{\small\color{#1}#3}}


% \newcommand{\cvreference}[7]{%
%     \textbf{#1}\newline% Name
%     \ifthenelse{\equal{#2}{}}{}{\addresssymbol~#2\newline}%
%     \ifthenelse{\equal{#3}{}}{}{#3\newline}%
%     \ifthenelse{\equal{#4}{}}{}{#4\newline}%
%     \ifthenelse{\equal{#5}{}}{}{#5\newline}%
%     \ifthenelse{\equal{#6}{}}{}{\emailsymbol~\texttt{#6}\newline}%
%     \ifthenelse{\equal{#7}{}}{}{\phonesymbol~#7}}

\begin{document}

\makecvtitle % Print the CV title
\vspace{-13mm}
\section{Introduction}
\cvitem{Goal: } {Seeking an architect or a principal leadership role in the area of network infrastructure and technologies, leveraging my experience in research, product \& standards development.}
\cvitem{Industry Leadership: } {I am a data- and tele-communication industry  verteran, having developed carrier-grade software for ISP markets (MPLS, VPLS, QoS, Policy based routing). I have a deep knowledge of different types of network architectures and protocols (e.g. ISP, data-center, mobile-networks, home-networks, industrial networks, IoT networks, etc.). Involved in next-gen low-latency, application-aware, data plane technologies for 5G, 6G, industrial networks) and have a working knowledge of toolsets XDP, eBPF, P4 etc. I am known for successful collaborations with internal and external partners, hold leadership roles in standards (IETF, ITU, ETSI) and research community (IEEE, ACM). \newline I have a number of patents my work is peer-reviewed \& published in both academic and industry standard organizations.}
%\cvitem{Proven Record: } {I have been in the data- and tele-communication industry  throughout my career and have a track record of accomplishing beyond the project goals. Known for successful collaborations with internal and external partners, active in standards (IETF, ITU, ETSI) and research community (IEEE, ACM). \newline As a technologist, I have a number of patents, my work is peer-reviewed \& published in both academic and industry standard organizations.}
\cvitem{Strategy \& Planning: }{Reporting into the OCTO, I contribute to long-term technology roadmap, identifying opprtunities and timing; collaborated with executives and managers to architect from incubated ideas to realization, strategic initiatives, defining  targets, requirements, and actionable recommendations.}

 \cvitem{Leadership | Mentor: } {Guided and trained graduate students through industry proof-of-concept projects; mentored students at my alma mater during covid times on career, development projects, and soft skills as part of a level-up program.}

% {}

%----------------------------------------------------------------------------------------
%	WORK EXPERIENCE SECTION
%----------------------------------------------------------------------------------------

\section{Professional Experience}
\subsection{Futurewei Technologies, Santa Clara, CA}
\cventry{Nov, 2013 -- present}{\textit{Principal Engineer}}{}{}{}
{Working with the office of the CTO, involved in the industry research on next-generation (5G/B5G/6G) network technologies and protocols in service provider and large-scale network markets. Day to day work involves equal parts research, standardization and collaboration. It entails academic and public presentations on topics, meeting partners and customers explaining unique value propositions from our research; solution feasibility through technical design and PoCs.\\
}
\cventry{Cloudcasting}{\textit{A scalable virtual routing architecture for cloud centric applications}}{}{}
{Developed an overlay routing architecture and multi-protocol dataplane auto-provisioning of tenant networks, specifically for cloud centric networks}{}
\cventry{5G/B5G dataplanes}{\textit{Innovations in IP Networks}}{}{}
{Focus on \textit{in-network services} and capabilities for emerging applications in industry automation, and advanced media ('holographic type' and high-precision communications). Analyse challenges in current networks and research new data plane and packet format mechanisms. Designing proof-of-concepts for research in transport and network layer}{}
\cventry{Network Slicing}{\textit{Operator-centric Architecture for Network Slices}}{}
{5G/B5G Network Slicing configuration, on-boarding, and  distributed protocol for resource sharing in a multi-domain, multi-provider network}{Early proponent of standardizing Network slice framework} {}
\cventry{5G Transport}{\textit{3GPP Mobile Backhaul Transport Networks}}{}{}
{An in-depth study and evaluation of end-to-end LTE, 5G/B5G Network Slicing architecture with specific focus on mobility and protocol overheads in mobile networks} {}

\cventry{Leadership in Standards:} {\textit{ITU FG-NET-2030}}{}{}
 {\emph{is a focus group on network technologies for 2030 and beyond in ITU}. Was instrumental to the establishment and conception of the FG. 
 %To develop formal road map and requirements for the next-gen network technologies
 } {At IETF, developing standards relating to 5G/B5G networks in IETF, specifically in IoT and DETNETs.}

%\cventry{Leadership: } {\textit{Project Management}}{}{}
%{Long and short term strategy and planning for our projects with emphasis on adapting to  changing market conditions}
\cventry{Management}{\textit{University Collaborations}}{}{}
{collaborated with professors on next-gen research topics; reviewed, analysed, partenered in research activities.}{}

\subsection{Cisco Systems, San Jose, CA}
\cventry{July,2005 -- Oct,2013}{\textit{Technical Leader}}{}{}{}{Led design \& development of Edge routing and switching component on IOS-XR, IOS. Including first in-house support of an Openflow agent, policy based routing, Lawful intercept, etc. Solved high-availability  and scalability challenges on a distributed  system such as IOS-XR. \newline At Cisco, I have mentored new hires, resolved customer issues and worked with  platform-hardware and core operating system teams.}

\cvitem{Policy services layer :}{Developed architecture to support flow-based infrastructure as well as network programmability (ONEP) for a distributed system.  \textbf{Platform Stability}: Provided HA solution for lawful intercept as a part of L2L3 services team. Developed hybrid architecture endian aware message translation methods.
}
 \cvitem{Edge Router Services:} {Developed platform-dependent slow-path forwarding module, added performance tweaks for QoS, selectively drop packets, prioritizing control packets, and bug fixes.
 }

 \cvitem{L2 functions:} {Designed support for CFM 802.1ag standard over EoMPLS and different Metro Ethernet services. The services included CFM over VFI interfaces.
 }
\cventry{Pre-2005}{\textit{Software Engineer}}{}{}{Contributed to  network device drivers and RTOS projects at HCL Tech, Futuresoft, Alcatel, Tasman Networks}{}

%----------------------------------------------------------------------------------------
%	EDUCATION SECTION
%----------------------------------------------------------------------------------------

\section{Education}

\cventry{1994--1997 :}{Master of Computer Applications}{Computer Science Department }{Pune University)}{}{}
\cventry{1991--1994 :}{Bachelor of Science}{Delhi University}{Delhi, India}{}{}

%----------------------------------------------------------------------------------------
%	COMPUTER SKILLS SECTION
%----------------------------------------------------------------------------------------

\section{Technical skills}

\cvitem{Programming}{\emph{Experienced:} C, C++, \emph{Familiar:} Python, Bash, Java}
\cvitem{Network Technologies}{Network Architecture, Protocols, 5G Arch, 5G-Backhaul, Edge computing, Network Slices, Service Provider and Data Center Networks }
\cvitem{Network Protocols}{\emph{Developed:} STP,  CFM, BFD, QoS, PBR, lawful-intercept, {\emph{Fluency:} BGP, IS-IS, Segment Routing, and MPLS}, IIoT and IoT protocols.}
%\cvitem{Standards and Research }{Authored specifications for IETF, ITU-T, and ETSI. TPC, \textit{author, chair, reviewer} in ACM, IEEE involvement }
\cvitem{Platforms} {IOS, IOS-XR, ASR9k, 7600 series, XDP, Linux networking - namespaces}

\section{External Positions of Responsibility}

\cvitem{2021-  }{\textbf{IETF WG Chair} for the SNAC-WG: Stub Networks using IPv6.}
%\cvitem{2018}{\textbf{Session Chair} of the \textbf{$1^{st}$} New Internet Forum,  ETSI Track, at SDN NFV Congress}
\cvitem{2017-23}{\textbf{General Chair} of the Workshops at IEEE ICNP: NIPAA'20, NIPAA'21, NIB'22
and workshops at ACM SIGCOMM NEAT'18, NEAT'19 and IIoT-NETS'23.}
\cvitem{2017}{\textbf{ETSI Rapporteur} for Network Slicing Architecture, Next-generation Protocols group (NGP).}
\cvitem{2019, 21}{\textbf{Invited delegate} to annual ITU's CXO Advisory Board meetings for Telecommunication Standardization Bureau (ITU-TSB).}

\vspace{5mm}
\begin{comment}
\begin{center}
\begin{tikzpicture}
\begin{scope}[curve in scope]
	\path (0,0) -- node[auto] { End CV. (Published work listed below)} (10,0);
\end{scope}
 \end{tikzpicture}
 \end{center}

\section{Publications}

\subsection{Published Papers (Accepted)}
\newbibliography{journal}
\nocite{journal}{*}
\bibliographystyle{journal}{plainyrrev}
\bibliography{journal}{journal}
{\large \textsc{Refereed Journal Articles}}

\subsection{Standards and Industry Relevant Publications}
\newbibliography{conference}
\nocite{conference}{*}
\bibliographystyle{conference}{plainyrrev}
\bibliography{conference}{conference}
{\large \textsc{Standards and Industry Relevant Publications}}

\subsection{Publications and Speeches as Subject Matter Expert}
\newbibliography{speeches}
\nocite{speeches}{*}
\bibliographystyle{speeches}{plainyrrev}
\bibliography{speeches}{speeches}
{\large \textsc{Publications and Speeches as Subject Matter Expert}}

\end{comment}
\section{References}
Available on Request
\\
Public Profiles: see embedded links \textcolor[RGB]{110, 150, 240}{\href{https://scholar.google.com/citations?hl=en&user=whVRYl4AAAAJ}{ 'Google Scholar'} 
 and \href{https://datatracker.ietf.org/person/kiran.ietf@gmail.com}{'IETF person'}}
\end{document}
